\begin{theorem}[Markov’s Inequality]
Let \((\Omega,\mathcal{F},\mathbb{P})\) be a probability space.  
Assume \(X:\Omega\to[0,\infty]\) is a non‑negative random variable and let \(a>0\).  
Then
\[
\mathbb{P}(X\ge a)\le\frac{\mathbb{E}[X]}{a}.
\]
\end{theorem}

\begin{proof}
\begin{enumerate}
\item \textbf{Trivial cases.}  
If \(\mathbb{E}[X]=+\infty\) the right‑hand side is \(+\infty\) and the inequality is automatically true.  
If \(\mathbb{P}(X\ge a)=0\) the inequality reads \(0\le\mathbb{E}[X]/a\), which holds because \(\mathbb{E}[X]\ge0\).  
Hence we may assume \(0<\mathbb{P}(X\ge a)<1\) and \(\mathbb{E}[X]<\infty\).

\item \textbf{Indicator bound.}  
Define the indicator of the event \(\{X\ge a\}\) by
\[
I_{\{X\ge a\}}(\omega)=
\begin{cases}
1, & X(\omega)\ge a,\\[2pt]
0, & X(\omega)< a .
\end{cases}
\]
For each \(\omega\in\Omega\),
\begin{equation}
X(\omega)\ge a\,I_{\{X\ge a\}}(\omega). \tag{2.1}
\end{equation}
Verification of (2.1): if \(X(\omega)\ge a\) then \(I_{\{X\ge a\}}(\omega)=1\) and \(a\le X(\omega)\);  
if \(X(\omega)<a\) then \(I_{\{X\ge a\}}(\omega)=0\) and (2.1) reduces to \(X(\omega)\ge0\), true since \(X\ge0\) a.s.

\item \textbf{Expectation of the bound.}  
Both sides of (2.1) are non‑negative, so monotonicity of expectation yields
\begin{equation}
\mathbb{E}[X]\ge a\,\mathbb{E}\!\big[I_{\{X\ge a\}}\big]. \tag{3.1}
\end{equation}

\item \textbf{Expectation of an indicator.}  
By definition of expectation for indicator functions,
\begin{equation}
\mathbb{E}\!\big[I_{\{X\ge a\}}\big]=\mathbb{P}(X\ge a). \tag{4.1}
\end{equation}

\item \textbf{Combine (3.1) and (4.1).}  
Substituting (4.1) into (3.1) gives
\[
\mathbb{E}[X]\ge a\,\mathbb{P}(X\ge a),
\]
which is equivalent to
\begin{equation}
\mathbb{P}(X\ge a)\le\frac{\mathbb{E}[X]}{a}. \tag{5.1}
\end{equation}
This completes the proof for finite expectations; the case \(\mathbb{E}[X]=+\infty\) was settled in step 1.

\item \textbf{Equality condition.}  
Equality in (5.1) can occur only if equality holds both in the pointwise bound (2.1) and in the expectation monotonicity step (3.1).

\begin{itemize}
\item \emph{Pointwise bound:} equality in (2.1) requires
\[
X(\omega)=a\,I_{\{X\ge a\}}(\omega)\quad\text{for }\mathbb{P}\text{-a.e. }\omega,
\]
i.e.
\[
X(\omega)=
\begin{cases}
a, & X(\omega)\ge a,\\
0, & X(\omega)< a .
\end{cases}
\]

\item \emph{Expectation step:} if the above identity holds \(\mathbb{P}\)-a.s., then taking expectations gives
\[
\mathbb{E}[X]=a\,\mathbb{P}(X\ge a),
\]
so (3.1) and consequently (5.1) become equalities.
\end{itemize}
Thus
\[
\mathbb{P}(X\ge a)=\frac{\mathbb{E}[X]}{a}
\iff
X=a\,I_{\{X\ge a\}}\quad\text{a.s.},
\]
i.e. \(X\) assumes only the two values \(0\) and \(a\) (up to a \(\mathbb{P}\)-null set). In all other situations the inequality is strict.
\end{enumerate}
\end{proof}